\section*{Capítulo \ref{ch:b2:microbial-metabolism} --- Metabolismo microbiano e biofilmes}

\begin{solution}{exo:b2:microbial-metabolism:1}
Cianobactéria: fotolitoautotrófica (luz, água, $\mathrm{CO_2}$).
\emph{Nitrosomonas}: quimiolitoautotrófica (oxidação do amônio,
$\mathrm{CO_2}$). \emph{E.~coli} em glicose: quimiorganoheterotrófica.
Bactéria púrpura não sulfurosa sobre succinato na luz:
fotorganoheterotrófica. \emph{Thiobacillus} sobre sulfeto no escuro:
quimiolitoautotrófica.
\end{solution}

\begin{solution}{exo:b2:microbial-metabolism:2}
Oxigênio (os milímetros da superfície), nitrato (logo abaixo, onde o
oxigênio acabou), óxidos de manganês e de ferro (a transição do marrom
para o cinza), sulfato (a camada negra de sulfeto, espessa nos
sedimentos marinhos), $\mathrm{CO_2}$ para a \omterm{def:b2:microbial-metabolism:anaerobic}{metanogênese} (o mais
profundo, ou logo abaixo da superfície na água doce pobre em sulfato).
\end{solution}

\begin{solution}{exo:b2:microbial-metabolism:3}
$5\,\mathrm{C_6H_{12}O_6} + 24\,\mathrm{NO_3^-} + 24\,\mathrm{H^+}
\to 30\,\mathrm{CO_2} + 12\,\mathrm{N_2} + 42\,\mathrm{H_2O}$: 4.8
íons nitrato por glicose (cada nitrato aceita 5 elétrons, cada glicose
fornece 24).
\end{solution}

\begin{solution}{exo:b2:microbial-metabolism:4}
Uma comunidade microbiana presa a uma superfície e imersa numa matriz
de polissacarídeos, proteínas e DNA que ela mesma produz. Etapas:
adesão, microcolônia e secreção da matriz, maturação com torres e
canais, dispersão. A placa bacteriana e as cáries; as infecções de
cateteres e implantes (e dos pulmões na fibrose cística); a
incrustação e a corrosão de tubulações e cascos --- ou, de modo útil,
os filtros biológicos percoladores do tratamento de água.
\end{solution}

\begin{solution}{exo:b2:microbial-metabolism:5}
Oito elétrons caem de \qty{-0.42}{V} para \qty{-0.24}{V}:
$\Delta G^{\circ\prime} = -8\times 96.5\times 0.18 =
\qty{-139}{kJ/mol}$ de metano (valor tabelado \qty{-131}{kJ/mol}), isto
é, \qty{35}{kJ} por $\mathrm{H_2}$: no máximo 0.7 ATP por hidrogênio,
na prática cerca de meio. A energia por molécula de substrato é tão
pequena que os metanogênicos crescem devagar (\omterm{def:b2:unicellular-diversity:division}{tempos de geração} de
horas a dias) e convertem quase todo o seu substrato em gás, e não em
células.
\end{solution}

\begin{solution}{exo:b2:microbial-metabolism:6}
$\mu = 1.2\,S/(5 + S)$: \qty{0.20}{h^{-1}} ($T = \ln 2/\mu =
\qty{3.5}{h}$), \qty{0.60}{h^{-1}} (\qty{1.2}{h}), \qty{1.09}{h^{-1}}
(\qty{0.64}{h}). Noventa por cento: $S = 9K_s = \qty{45}{mg/L}$.
\end{solution}

\begin{solution}{exo:b2:microbial-metabolism:7}
$S^* = 0.05\times 0.4/(0.8 - 0.4) = \qty{0.05}{g/L}$; $X^* =
0.4\times(5 - 0.05) = \qty{1.98}{g/L}$; produção $DX^* = 0.4\times
1.98 = \qty{0.79}{g/L/h}$; lavagem $D_c = 0.8\times 5/5.05 =
\qty{0.79}{h^{-1}}$.
\end{solution}

\begin{solution}{exo:b2:microbial-metabolism:8}
Capturado: $0.25\times 275 = \qty{69}{kJ}$, isto é, 1.4 ATP por
amônio; $12/1.4 = 8.7$, cerca de nove íons amônio por carbono --- e
várias vezes mais quando se conta o custo de produzir NADH empurrando
elétrons ladeira acima a partir do nitrito. Um heterótrofo recebe seu
carbono já reduzido e organizado, e trinta ATP por glicose; uma
nitrificante precisa comprar tanto sua energia quanto seu poder
redutor de um doador pobre e construir cada carbono a partir do
$\mathrm{CO_2}$, de modo que converte em células uma fração ínfima do
substrato que processa.
\end{solution}

\begin{solution}{exo:b2:microbial-metabolism:9}
$L_p = \sqrt{2\times 1.2\times 10^{-9}\times 0.2/0.5} =
\sqrt{9.6\times 10^{-10}} = \qty{31}{\micro m}$. Dobrar $c_0$
multiplica $L_p$ por $\sqrt 2$: \qty{44}{\micro m}; quadruplicar $q$
o reduz à metade: \qty{15}{\micro m}.
\end{solution}

\begin{solution}{exo:b2:microbial-metabolism:10}
De cima para baixo: água com algas e cianobactérias (fotossíntese
oxigênica); uma faixa cor de ferrugem de oxidantes de ferro, onde o
$\mathrm{Fe^{2+}}$ encontra o oxigênio; um véu branco de oxidantes de
enxofre incolores (\omterm{def:b2:microbial-metabolism:trophic}{quimiolitotróficos}), onde o sulfeto encontra o
oxigênio; uma faixa púrpura de bactérias sulfurosas púrpuras
(\omterm{def:b2:microbial-metabolism:anoxygenic}{fotossíntese anoxigênica} sobre sulfeto); uma faixa verde de bactérias
sulfurosas verdes (o mesmo, tolerando mais sulfeto e menos luz); lama
negra de fermentadores e redutores de sulfato.
(a) As bactérias sulfurosas púrpuras precisam ao mesmo tempo da luz
que vem de cima e do sulfeto que vem de baixo, por isso ficam na
interface onde os dois gradientes se cruzam. (b) Os redutores de
sulfato produzem $\mathrm{H_2S}$, que precipita o ferro como FeS
negro. (c) Fechada para a matéria: o enxofre, o carbono e o nitrogênio
circulam entre formas oxidadas e reduzidas sem sair; aberta para a
energia: a luz entra, o calor sai, e sem luz todos os ciclos param.
\end{solution}

\begin{solution}{exo:b2:microbial-metabolism:11}
Produção $pn$, perda $kA$: $A^* = pn/k$. $A_c = 10^{-8}\times
6.02\times 10^{23} = 6.0\times 10^{15}$ moléculas por litro; $n_c =
kA_c/p = 5\times 10^{-4}\times 6.0\times 10^{15}/5 = 6\times 10^{11}$
por litro, $6\times 10^{8}$ por mililitro. A cultura densa
($10^{9}$) está acima do limiar; a água do mar ($10^{5}$), quatro
ordens de grandeza abaixo. Num \omterm{def:b2:microbial-metabolism:biofilm}{biofilme}, as células ficam a $10^{11}$
por mililitro de matriz, e a matriz retarda a saída do sinal
(o que equivale a reduzir $k$), de modo que uma microcolônia de alguns
milhares de células num nanolitro atinge o quórum, ao passo que as
mesmas células dispersas num litro nunca o atingiriam.
\end{solution}

\begin{solution}{exo:b2:microbial-metabolism:12}
Fixação de nitrogênio (só os \omterm{def:b2:unicellular-diversity:prokaryote}{procariotos} têm nitrogenase): se ela
parasse, o nitrogênio combinado da biosfera escoaria por
\omterm{def:b2:microbial-metabolism:anaerobic}{desnitrificação} e a produção entraria em colapso em poucas décadas.
Nitrificação e \omterm{def:b2:microbial-metabolism:anaerobic}{desnitrificação} (apenas \omterm{def:b2:unicellular-diversity:prokaryote}{procariotos}): se parassem, o
amônio se acumularia, o nitrato desapareceria e o nitrogênio nunca
voltaria ao ar. \omterm{def:b2:microbial-metabolism:anaerobic}{Metanogênese} e \omterm{def:b2:microbial-metabolism:anaerobic}{respiração anaeróbia} (apenas
\omterm{def:b2:unicellular-diversity:prokaryote}{procariotos}): se parassem, a matéria orgânica dos sedimentos
anóxicos, dos rúmens e dos pântanos se acumularia sem ser
mineralizada e o carbono sairia do ciclo; do mesmo modo, a fotossíntese
oxigênica foi uma invenção bacteriana, e sem ela não haveria oxigênio
para respirar. Os eucariotos são um acréscimo tardio e quimicamente
estreito a um planeta procariótico.
\end{solution}

\begin{solution}{pb:b2:microbial-metabolism:1}
\textbf{1.} $-24\times 96.5\times (0.82 + 0.43) = \qty{-2900}{kJ/mol}$.
\textbf{2.} Nitrato: $-24\times 96.5\times 1.17 = \qty{-2700}{kJ/mol}$;
sulfato: $-24\times 96.5\times 0.21 = \qty{-490}{kJ/mol}$.
\textbf{3.} O oxigênio rende mais, de modo que os aeróbios superam os
desnitrificantes na competição pela matéria orgânica enquanto houver
oxigênio, e o oxigênio reprime as nitrato-redutases; a \omterm{def:b2:microbial-metabolism:anaerobic}{redução do
sulfato} exige anoxia \emph{e} ausência de nitrato, de modo que seu
cheiro indica que o tanque está subaerado e o nitrato esgotado --- e o
$\mathrm{H_2S}$ é tóxico e corrói o concreto.
\textbf{4.} $0.4\times 2900/50 = 23$ ATP; $0.4\times 2700/50 = 22$;
$0.4\times 490/50 = 4$.
\textbf{5.} Cerca de seis vezes mais ($23/4$).
\textbf{6.} Dois ATP por glicose significam um rendimento de cerca de
\num{0.1}, de modo que \qty{90}{\%} da energia do substrato permanece
nos produtos (ácidos, álcoois, $\mathrm{H_2}$), que os metanogênicos
convertem em metano com um ganho igualmente pequeno; pouca energia,
pouca biomassa, e o carbono sai como gás.
\textbf{7.} $10^4\times 0.3 = \qty{3000}{kg}$ de DQO por dia.
\textbf{8.} Biomassa $0.4\times 3000 = \qty{1200}{kg}$ de DQO; os
\qty{1800}{kg} de DQO restantes são oxidados e exigem \qty{1800}{kg}
de $\mathrm{O_2}$.
\textbf{9.} $1800/2 = \qty{900}{kWh}$ por dia.
\textbf{10.} $10^4\times 0.04 = \qty{400}{kg}$ de nitrogênio;
$4.57\times 400 = \qty{1830}{kg}$ de $\mathrm{O_2}$.
\textbf{11.} Em regime estacionário, $\mu = 1/\theta \le \mu_{\max}$:
lavagem abaixo de $\theta = 1/\mu_{\max} = \qty{2}{d}$.
\textbf{12.} $D = \qty{0.1}{d^{-1}}$: $S^* = 1\times 0.1/(0.5 - 0.1)
= \qty{0.25}{g/m^3}$.
\textbf{13.} $S^* = 0.1/(0.25 - 0.1) = \qty{0.67}{g/m^3}$, e o limite
de lavagem sobe para \qty{4}{d}; o operador precisa aumentar a idade
do lodo (descartar menos lodo): com $\theta = \qty{20}{d}$, $S^* =
0.05/0.2 = \qty{0.25}{g/m^3}$ de novo.
\textbf{14.} $2.86\times 400 = \qty{1144}{kg}$ de DQO por dia.
\textbf{15.} Amônio no efluente $0.25\times 10^4 = \qty{2.5}{kg/d}$,
\qty{0.6}{\%} da carga; o restante, \qty{397}{kg} de nitrogênio, sai
como $\mathrm{N_2}$: \qty{99}{\%}.
\textbf{16.} Economizados: \qty{1144}{kg} de $\mathrm{O_2}$. Demanda
total: $(1800 - 1144) + 1830 = \qty{2490}{kg}$ de $\mathrm{O_2}$ por
dia.
\textbf{17.} $1200\times 0.6\times 0.35 = \qty{252}{m^3}$ de metano
por dia.
\textbf{18.} $252\times 36 = \qty{9070}{MJ} = \qty{2520}{kWh}$;
eletricidade $0.35\times 2520 = \qty{880}{kWh}$ por dia.
\textbf{19.} Aeração $2490/2 = \qty{1245}{kWh}$ por dia: o metano
cobre $880/1245 = \qty{71}{\%}$ dela.
\textbf{20.} $D_e\,c'' = q$, logo $c = c_0 + ax + qx^2/2D_e$; com
$c(L_p) = 0$ e $c'(L_p) = 0$: $a = -qL_p/D_e$ e $c =
(q/2D_e)(L_p - x)^2$, com $L_p^2 = 2D_ec_0/q$.
\textbf{21.} $L_p = \sqrt{2\times 1.5\times 10^{-9}\times 0.2/0.3} =
\sqrt{2\times 10^{-9}} = \qty{45}{\micro m}$.
\textbf{22.} $45/300 = \qty{15}{\%}$ aeróbia. Abaixo, os
desnitrificantes reduzem a $\mathrm{N_2}$ o nitrato que se difunde a
partir da camada nitrificante da superfície, usando a matéria orgânica
que se difunde para dentro.
\textbf{23.} Num mesmo filme, o gradiente de oxigênio cria uma camada
aeróbia (nitrificação) sobre uma anóxica (\omterm{def:b2:microbial-metabolism:anaerobic}{desnitrificação}), a algumas
dezenas de micrômetros de distância; o tanque aerado é óxico em toda a
sua extensão, de modo que a \omterm{def:b2:microbial-metabolism:anaerobic}{desnitrificação} exige um tanque anóxico
separado.
\textbf{24.} O cloro reage com os polímeros da matriz e é consumido
por eles antes de alcançar as células profundas; as células profundas
estão famintas, crescem devagar ou estão dormentes e são muito menos
sensíveis; só a camada superficial morre, e o filme volta a crescer a
partir de baixo.
\textbf{25.} Demanda de oxigênio de \qty{2490}{kg} por dia; metano,
\qty{252}{m^3} por dia; o metano fornece \qty{71}{\%} da eletricidade
de aeração (\qty{880}{kWh} de \qty{1245}{kWh}).
\end{solution}
